\section*{Capítulo \ref{ch:g6:matter-from-food} --- Matéria a partir do alimento}

\begin{solution}{exo:g6:matter-from-food:1}
Os \omterm{def:g5:food-webs:roles}{consumidores} transformam em \omterm{def:g6:plants-are-producers:matter}{matéria orgânica} própria a \omterm{def:g6:plants-are-producers:matter}{matéria
orgânica} que comem. As produtoras fabricam \omterm{def:g6:plants-are-producers:matter}{matéria orgânica} nova a
partir de suprimentos puramente minerais, com a \omterm{prop:g3:balanced-diet:jobs}{energia} da \omterm{def:g6:exploring-our-environment:conditions}{luz}.
\end{solution}

\begin{solution}{exo:g6:matter-from-food:2}
A \omterm{def:g4:journey-of-food:digestion}{digestão} desmonta o capim em \omterm{def:g4:journey-of-food:digestion}{nutrientes}
(\cref{ch:g4:journey-of-food}); o \omterm{prop:g4:journey-of-food:absorb}{sangue} entrega esses \omterm{def:g4:journey-of-food:digestion}{nutrientes} em
toda parte (\cref{ch:g4:heart-and-blood}); as \omterm{def:g5:first-look-at-cells:cell}{células} do corpo os
remontam em matéria de vaca, dividindo-se conforme o corpo cresce
(\cref{ch:g5:first-look-at-cells}).
\end{solution}

\begin{solution}{exo:g6:matter-from-food:3}
Crescimento (os quilos que o bezerro acrescenta); renovações (pelo
que cai, \omterm{def:g1:animals-around-us:covering}{penas} novas, chifres); reparos (um \omterm{def:g3:bones-and-muscles:skeleton}{osso} consertado); produtos
(\omterm{def:g2:animals-grow-up:born}{leite}, \omterm{def:g2:animals-grow-up:born}{ovos}, lã).
\end{solution}

\begin{solution}{exo:g6:matter-from-food:4}
Nunca absorvida --- sai como fezes; queimada com o oxigênio da
\omterm{def:g4:breathing:movements}{respiração} para movimento e calor --- sai como gás expirado; e a
sobra, construída em crescimento e produtos.
\end{solution}

\begin{solution}{exo:g6:matter-from-food:5}
Porque os \qty{25}{kg} de \omterm{def:g2:animals-grow-up:born}{leite} de cada dia são \omterm{def:g6:plants-are-producers:matter}{matéria orgânica} que
ela precisa reconstruir a partir da ração daquele dia; o \omterm{def:g2:animals-grow-up:born}{leite} é
produção, e produção é alimento reconstruído --- sem entrada, sem
saída.
\end{solution}

\begin{solution}{exo:g6:matter-from-food:6}
Montada \emph{dentro} da galinha, a partir dos \omterm{def:g4:journey-of-food:digestion}{nutrientes} do grão
dela --- mas a própria \omterm{def:g6:plants-are-producers:matter}{matéria orgânica} foi fabricada antes dela, nas
\omterm{def:g1:plants-around-us:parts}{folhas} da \omterm{def:g1:plants-around-us:plant}{planta} do grão. A galinha transformou; a \omterm{def:g1:plants-around-us:plant}{planta} produziu.
\end{solution}

\begin{solution}{exo:g6:matter-from-food:7}
Em cada seta, quase toda a matéria comida é queimada para viver ou
perdida nas fezes, e só uma fração modesta vira corpo de quem comeu
--- a única parte que o andar seguinte pode comer. Cada andar tem,
portanto, só uma fração da matéria do andar de baixo: a pirâmide
estreita por aritmética.
\end{solution}

\begin{solution}{exo:g6:matter-from-food:8}
Entra: carne ou ração (e um camundongo de vez em quando). Construído:
o crescimento dele quando gatinho, o pelo que ele renova a cada muda,
os reparos depois de arranhões. Sai: fezes no quintal e a fatia
queimada que dá força a cada bote, a cada ronronar e a cada soneca
quentinha.
\end{solution}

\begin{solution}{exo:g6:matter-from-food:9}
O que difere é a fonte e o equipamento: capim contra coelho,
trituradores com intestino enorme contra agarradores com intestino
curto. O que é idêntico é o próprio ofício: os dois desmontam a
\omterm{def:g6:plants-are-producers:matter}{matéria orgânica} comida em \omterm{def:g4:journey-of-food:digestion}{nutrientes} e a reconstroem como substância
própria.
\end{solution}

\begin{solution}{exo:g6:matter-from-food:10}
Quase tudo saiu como a fatia queimada --- dando força a um ano de
correr, pensar e ficar a \qty{37}{\degreeCelsius} --- e uma parte
nunca foi absorvida e saiu como fezes. Só a sobra pequena, os
\qty{3}{kg}, virou corpo novo: o balanço da criança fecha como o da
vaca.
\end{solution}

\begin{solution}{exo:g6:matter-from-food:11}
O corpo do caracol extraiu os ingredientes minerais da \omterm{def:g1:animals-around-us:covering}{concha} do
alimento e do que estava em volta dele, transportou esses
ingredientes e os depositou camada a camada na arquitetura dele. A
substância é mineral; a \emph{construção} é trabalho do \omterm{def:g1:animals-around-us:animal}{animal} vivo
--- uma coisa fabricada, com material mineral.
\end{solution}

\begin{solution}{exo:g6:matter-from-food:12}
Passando pelo porco, quase toda a matéria das batatas é gasta no
próprio viver do porco --- queimada para calor e movimento, perdida
nas fezes --- e só uma fração volta como carne. Comidas diretamente,
as batatas pulam as despesas do porco. A diferença foi para o balanço
do porco, e não para a mesa.
\end{solution}

\begin{solution}{pb:g6:matter-from-food:1}
\textbf{1.} Os capins do pasto e as \omterm{def:g1:plants-around-us:plant}{plantas} da horta (e qualquer
fruteira). Ingredientes: água e \omterm{prop:g4:what-plants-need:needs}{sais minerais} do solo, gás carbônico
do ar. \omterm{prop:g3:balanced-diet:jobs}{Energia}: a \omterm{def:g6:exploring-our-environment:conditions}{luz}.

\textbf{2.} Nas próprias \omterm{def:g1:plants-around-us:parts}{folhas} dos pés de capim --- cada grama: as
produtoras fabricam no lugar, de dia.

\textbf{3.} Devolver o que se toma emprestado: o esterco, ao
apodrecer, reabastece os \omterm{prop:g4:what-plants-need:needs}{sais minerais} do solo da horta, repondo o
que as colheitas levam embora.

\textbf{4.} Porque a vaca é uma consumidora: o ofício dela precisa de
matéria \emph{orgânica} para desmontar e reconstruir. \omterm{prop:g4:what-plants-need:needs}{Sais minerais} e
água são ingredientes de produtora; ela não tem maquinaria para
construir corpo a partir de minerais --- é justamente esse o ofício
que lhe falta.

\textbf{5.} Uma parte passa sem ser absorvida --- as fezes; uma parte
grande é queimada com o oxigênio respirado para dar força a andar,
mastigar e se manter quente; o resto vira \omterm{def:g2:animals-grow-up:born}{leite} e vaca (e o balanço
do dia fecha).

\textbf{6.} Entra: grão. Construído: \omterm{def:g2:animals-grow-up:born}{ovos}, a manutenção dela e ---
na muda --- uma plumagem nova inteira. Sai: fezes e a fatia queimada.
Na muda, a linha do ``construído'' é gasta com \omterm{def:g1:animals-around-us:covering}{penas}; com a entrada
igual, a rubrica dos \omterm{def:g2:animals-grow-up:born}{ovos} precisa encolher durante um tempo.

\textbf{7.} Capim $\rightarrow$ galinha (via grão, que já é um produto
vegetal) $\rightarrow$ pintinho $\rightarrow$ gavião; ou, do lado do
pasto: capim $\rightarrow$ vaca. A matéria foi \emph{fabricada} só nos
elos verdes (capim, \omterm{def:g1:plants-around-us:plant}{plantas} do grão); todo elo \omterm{def:g1:animals-around-us:animal}{animal} --- galinha,
pintinho, gavião, vaca --- apenas a transformou.

\textbf{8.} Pão: \omterm{def:g1:plants-around-us:parts}{folhas} do pé de trigo. \omterm{def:g2:animals-grow-up:born}{Ovo} cozido: a galinha, a
partir do grão --- \omterm{def:g1:plants-around-us:parts}{folhas} da \omterm{def:g1:plants-around-us:plant}{planta} do grão. \omterm{def:g2:animals-grow-up:born}{Leite}: a vaca, a partir
do capim --- \omterm{def:g1:plants-around-us:parts}{folhas} do capim. Três trilhas, três \omterm{def:g1:plants-around-us:parts}{folhas}, todas na \omterm{def:g6:exploring-our-environment:conditions}{luz}.

\textbf{9.} Cada andar repassa só uma fração: dos \qty{10000}{kg} de
capim, uma população de coelhos constrói como corpo apenas a fração
modesta dela; dessa fração, as raposas constroem de novo só uma
fração. Dois andares de perdas encolhem o capim de um pasto até uma
tocazinha de raposas --- a estimativa dela é a aritmética da pirâmide.

\textbf{10.} Comida diretamente, a produção da horta chega à família
sem nenhum andar de \omterm{def:g5:food-webs:roles}{consumidor} no meio: nada dela é queimado no viver
de um \omterm{def:g1:animals-around-us:animal}{animal} nem perdido nas fezes dele. A rota do pasto paga
primeiro todas as despesas da vaca, então cada metro quadrado
alimenta menos gente.

\textbf{11.} O capim para de produzir, então a entrada da vaca falha:
primeiro o \omterm{def:g2:animals-grow-up:born}{leite} encolhe (a produção é alimento reconstruído), depois
ela queima as próprias reservas de gordura e perde peso --- a linha do
``construído'' fica negativa. O sol é \omterm{prop:g3:balanced-diet:jobs}{energia}, e não substância: sem
água não há com o que o capim construir, então o pasto não pode ser
salvo pelo céu.

\textbf{12.} Produtoras: as \omterm{def:g1:plants-around-us:plant}{plantas} do sítio fabricam toda a \omterm{def:g6:plants-are-producers:matter}{matéria
orgânica} dele a partir de suprimentos minerais e da \omterm{def:g6:exploring-our-environment:conditions}{luz}.
\omterm{def:g5:food-webs:roles}{Consumidores}: a vaca, as galinhas e a família transformam essa
matéria, queimando quase tudo e construindo o resto. Faltando: os
\omterm{def:g5:food-webs:roles}{decompositores} --- que precisam devolver as fezes, as bostas de vaca e
os restos do outono à \omterm{def:g6:plants-are-producers:matter}{matéria mineral}; eles são o assunto do próximo
capítulo.
\end{solution}
